\cotefile{64}
\scanpage{116}
% archive sheet number: 115
soit
\[
 \sigma_i z=\frac{\unsure{\zeta_i|\xi_i}z+\overbrace{\zeta_j\zeta_k-\zeta_i(\zeta_j+\zeta_k)}^{b_i}}{z-\zeta_i}
 \tag*{(61)}
\]
On peut transformer $b_i$, en utilisant les fonctions symétriques élémentaires $(-1)^\alpha a_\alpha$ des $\zeta$,
\[
 b_i=2\underbrace{\zeta_j\zeta_k}_{-a_3/\zeta_i}-
 (\underbrace{\zeta_i\zeta_j+\zeta_j\zeta_k+\zeta_k\zeta_i}_{a_2})
 =-(a_2+2a_3/\zeta_i)
\]
donc on trouve
\[
 \sigma_i z=\frac{\zeta_i z-(a_2+2a_3/\zeta_i)}{z-\zeta_i}
 =\frac{\zeta_i^2z-(a_2\zeta_i+2a_3)}{\zeta_i z-\zeta_i^2}
 \tag*{(62)}
\]
% slanted
\margin{expression valable du moins si $\zeta_i\ne0$, sinon on convient bien sûr que $a_3/\zeta_i$ \deleted{\unsure{se}} devient $\zeta_j\zeta_k$. Mais quitte à modifier en ajoutant une cte, \unsure{ts} les $\zeta_i\ne0$ \ldots}
Les \inserted{div. $\Delta_i$ des} points fixes de $\sigma_i$ est l'\deleted{ens} div. des solutions de l'équation en $z_i$
\[
 z_i^2-2\zeta_i z_i+(a_2+2a_3/\zeta_i)=0 \tag*{(63)}
\]
dont les solutions sont données par
\[
 z_i=\zeta_i\pm\sqrt{\underbrace{\zeta_i^2-a_2-2a_3/\zeta_i}_{\deleted{\unsure{\zeta}}}} \tag*{(64)}
\]
L'expression sous le signe $\sqrt{\vphantom{x}}$ s'écrit aussi
\[
 \frac{1}{\zeta_i}(\zeta_i^3-a_2\zeta_i-2a_3)
\]
ou encore, compte tenu de (66)
\[
 \zeta_i^3=-a_1\zeta_i^2-a_2\zeta_i-a_3
\]
\deleted{ce radical} cette expression devient
\[
 -\frac{1}{\zeta_i}(a_1\zeta_i^2+2a_2\zeta_i+3a_3)
\]
donc (\unsure{64|69}) s'écrit aussi
\[
 \left\{\begin{aligned}
 (z_i-\zeta_i)^2&=-\frac{1}{\zeta_i}(a_1\zeta_i^2+2a_2\zeta_i+3a_3)\\
 \text{i.e}\quad z_i&=\zeta_i\pm\sqrt{-\tfrac1{\zeta_i}(a_1\zeta_i^2+2a_2\zeta_i+3a_3)}
 \end{aligned}\right. \tag*{\textcircled{65}}
\]
% slanted
\margin{Si $\zeta_i=0$, on \deleted{\unsure{effectuera division}} \inserted{écrit} $-(a_1\zeta_i+2a_2+3\zeta_j\zeta_k)$ l'expr. sous le radical.}
Les douze \unsure{point|points} d'ordre $4$ de $E$ sont donc les points au dessus des six points $z_i$, donc \deleted{définis par} \inserted{un tel}
